\lysh{16151_SOLUTION}{1}
ΛΥΣΗ
α) $\lambda_{\vec{\alpha}} = \frac{3}{3} = 1$ και $\lambda_{\vec{\beta}} = \frac{1}{-\sqrt{3}} = -\frac{\sqrt{3}}{3}$. Αν $\omega$ είναι η γωνία που σχηματίζει το διάνυσμα $\vec{\alpha}$ με τον άξονα $x'x$ και $\phi$ η γωνία που σχηματίζει το διάνυσμα $\vec{\beta}$ με τον άξονα $x'x$, τότε $\text{εφ}\omega=\lambda_{\vec{\alpha}}=1$ και $\text{εφ}\phi = \lambda_{\vec{\beta}} = -\frac{\sqrt{3}}{3}$. Το διάνυσμα $\vec{\alpha}$ βρίσκεται στο $1^\text{ο}$ τεταρτημόριο, άρα $\omega=45^\circ$ και το διάνυσμα $\vec{\beta}$ βρίσκεται στο $2^\text{ο}$ τεταρτημόριο, άρα $\phi=150^\circ$.

β) Η γωνία $(\widehat{\vec{\alpha},\vec{\beta}})$ ισούται με τη γωνία που σχηματίζει το διάνυσμα $\vec{\beta}$ με τον άξονα $x'x$ αν αφαιρέσουμε τη γωνία που σχηματίζει το διάνυσμα $\vec{\alpha}$ με τον $x'x$.
Δηλαδή $(\widehat{\vec{\alpha},\vec{\beta}}) = 150^\circ - 45^\circ = 105^\circ$.

% TODO: TikZ σχήμα (μην το κατασκευάσεις)
\end{lysh}